\section{Related Work}


\noindent \textbf{Brain Activity Recognition.}
Earlier studies~\cite{vanrullen2019reconstructing,dado2022hyperrealistic,shen2019deep,seeliger2018generative,NEURIPS2022_bee5125b,gu2022decoding,ferrante2022semantic} aimed to reconstruct visual stimuli by decoding brain signals into the latent space of Generative Adversarial Networks~\cite{goodfellow2020generative}. Given the significant capabilities of diffusion models in image generation~\cite{rombach2022high,ho2020denoising}, numerous investigations~\cite{takagi2023high,chen2023seeing,scotti2023reconstructing,lu2023minddiffuser,bai2023dreamdiffusion,lan2023seeing,zeng2023controllable,li2024neuraldiffuser,fang2024alleviating,li2024visual} have explored the use of their image prior for high-quality reconstruction. To improve reconstruction performance, some studies~\cite{zeng2023controllable,lu2023minddiffuser} exploit low-level signals within fMRI to achieve spatially consistent predictions by imposing spatial constraints during the generation process. Another set of studies~\cite{liu2023brainclip} has sought to leverage contrastive learning, such as CLIP, to improve semantic alignment in reconstruction. Certain works demonstrate the effectiveness of masked pretraining~\cite{he2022masked,bai2023dreamdiffusion} of brain signals on large datasets, facilitating robust latent space learning.
Moreover, novel applications related to brain activity~\cite{sun2023tuning,sun2024neurocine,wang2024mindbridge,nikolaus2024modality,xia2024umbrae,chen2024bridging,huo2024neuropictor} have emerged. For instance, Mind-Video~\cite{chen2023cinematic} extends image reconstruction to video reconstruction, while UniBrain~\cite{mai2023unibrain} not only reconstructs visual signals but also generates corresponding captions.

\paragraph{Image Manipulation.}
With a plethora of diverse generative priors embedded within powerful diffusion architectures, numerous studies~\cite{avrahami2023blended,couairon2022diffedit,meng2021sdedit,nichol2021glide,sun2023imagebrush,sun2024avi} have endeavored to adapt them for image manipulation. To facilitate the editing of desired semantic regions within spatial feature maps, researchers have proposed various techniques such as cross-attention injection~\cite{hertz2022prompt,balaji2022ediffi} and introduced semantic loss~\cite{tumanyan2023plug} to constrain spatial features during the generation process. Additionally, to provide explicit guidance towards intended instructions, subsequent works~\cite{brooks2023instructpix2pix,geng2023instructdiffusion} leverage language-based instructions to fine-tune the SD model, showcasing impressive capabilities in conforming to desired editing directions.
However, these approaches primarily operate on raw images. Few have delved into manipulating stored visual content with another modality, which holds potential for diverse applications such as fostering friendly human-computer interaction.

\paragraph{LLM Agent For Visual Assistance.}
Large language models (LLMs) have showcased profound capabilities as remarkable reasoning engines in various tasks~\citep{zhuge2023mindstorms,wei2022chain,huang2022towards,Touvron2023LLaMAOA,brown2020language}, due to their emerging ability~\citep{wei2022emergent}.
With rich contextual prior knowledge, LLMs adeptly tackle a multitude of visual-language tasks~\citep{zhuge2024language,alayrac2022flamingo,zhang2023multimodal,chakrabarty2023spy} via appropriate prompting adaptation.
Through visual instruction tuning, LLMs can discern image content, reason about involved events, and generate plausible responses~\citep{li2023blip,liu2024visual,zhu2023minigpt}. Subsequent studies~\citep{koh2024generating,sun2023generative,wu2023visual} have shown that LLMs can naturally provide visual feedback by leveraging an image rendering backbone such as a diffusion model.
Recently, researchers have employed LLMs to facilitate the image generation process~\citep{NEURIPS2023_3a7f9e48,fu2023guiding}, where the language model exhibits impressive capabilities in layout reasoning and instruction execution.
